\section*{अध्याय \ref{ch:b1:nucleophilic-substitution} --- नाभिकरागी प्रतिस्थापन}

\begin{solution}{exo:b1:nucleophilic-substitution:1}
\ce{CH3CH2OH + Br-}; \ce{CH3OCH2CH3 + I-}; \ce{CH3CH2CH2CN + Cl-};
\ce{CH3NH3+ + Br-}। \omterm{def:b1:nucleophilic-substitution:substitution}{क्रियाधार}: हैलोऐल्केन; \omterm{def:b1:electronic-effects:nucleophile}{नाभिकरागी}:
\ce{OH-}, \ce{CH3CH2O-}, \ce{CN-}, \ce{NH3}; \omterm{def:b1:electronic-effects:nucleophile}{निर्गामी समूह}: हैलाइड आयन।
\end{solution}

\begin{solution}{exo:b1:nucleophilic-substitution:2}
द्वितीय कोटि: तिगुनी; दोनों आधी करने पर 4 से भाग। केवल RY में प्रथम कोटि:
अपरिवर्तित; आधी।
\end{solution}

\begin{solution}{exo:b1:nucleophilic-substitution:3}
SN2 (प्राथमिक, अच्छा \omterm{def:b1:electronic-effects:nucleophile}{नाभिकरागी}, अप्रोटिक); SN1 (तृतीयक, जल); SN2
(मेथिल, प्रबल \omterm{def:b1:electronic-effects:nucleophile}{नाभिकरागी}); SN1 \omterm{def:b1:nucleophilic-substitution:sn1}{विलायक-अपघटन} (द्वितीयक, दुर्बल उदासीन
\omterm{def:b1:electronic-effects:nucleophile}{नाभिकरागी}, \omterm{def:b1:intermolecular-forces-solvents:solvent-classes}{प्रोटिक विलायक}), संभवतः कुछ SN2 के साथ।
\end{solution}

\begin{solution}{exo:b1:nucleophilic-substitution:4}
सायनाइड ब्रोमीन के विपरीत स्थान लेता है और तीनों अन्य
समूह पलट जाते हैं। उत्पाद में CN समूह पहले क्रम पर है, जैसे Br था,
अन्य क्रम अपरिवर्तित: प्रतीपित व्यवस्था का वर्णक
विपरीत है, \cip{S}-2-मेथिलब्यूटेननाइट्राइल।
\end{solution}

\begin{solution}{exo:b1:nucleophilic-substitution:5}
1. \omterm{def:b1:electronic-effects:cleavage}{विषमांश विदलन}: C--Br आबंध से Br की ओर तीर, जिससे \ce{(CH3)3C+} और
\ce{Br-} बनते हैं (मंद)। 2. जल का एक \omterm{def:b1:lewis-resonance-vsepr:lewis}{एकाकी युग्म} धनायनी कार्बन की ओर:
\ce{(CH3)3C-OH2+}। 3. जल के दूसरे अणु का एक \omterm{def:b1:lewis-resonance-vsepr:lewis}{एकाकी युग्म} \ce{OH2+} समूह का
एक प्रोटॉन ले लेता है, और O--H युग्म ऑक्सीजन पर लौट आता है। केवल पहला,
मंद चरण दर का निर्णय करता है; जल, जो विलायक है, इतने आधिक्य में है कि
उसकी सांद्रता नहीं बदलती।
\end{solution}

\begin{solution}{exo:b1:nucleophilic-substitution:6}
Br वहन करने वाला कार्बन प्राथमिक है, पर उसके बगल में एक टर्ट-ब्यूटिल समूह
C--Br आबंध के पीछे का स्थान भर देता है: SN2 \omterm{def:b1:elementary-steps:energy-profile}{संक्रमण अवस्था} में बहुत
भीड़ है। SN1 को एक प्राथमिक \omterm{def:b1:electronic-effects:intermediates}{कार्बधनायन} चाहिए होता, जो बहुत अस्थायी है।
\end{solution}

\begin{solution}{exo:b1:nucleophilic-substitution:7}
आयोडाइड SN2 द्वारा आयोडाइड का स्थान लेता है, हर बार प्रतीपन के साथ: \cip{S},
\cip{R} बनता जाता है जब तक दोनों समान रूप से उपस्थित न हों, यानी \omterm{def:b1:stereochemistry-in-depth:enantiomers}{रेसेमिक मिश्रण}। हर
प्रतिस्थापन एक \cip{S} अणु को एक \cip{R} में बदल देता है: घूर्णन
दो अणुओं के बराबर घटता है, इसलिए घूर्णन उस दर से दुगुनी तेज़ी से घटता है जिस दर से
\omterm{def:b1:nucleophilic-substitution:substitution}{क्रियाधार} प्रतिस्थापित होता है।
\end{solution}

\begin{solution}{exo:b1:nucleophilic-substitution:8}
SN2: ब्रोमोमेथेन $>$ 1-ब्रोमोप्रोपेन $>$ 2-ब्रोमोप्रोपेन $>$
2-ब्रोमो-2-मेथिलप्रोपेन (भीड़)। SN1: उलटा क्रम (\omterm{def:b1:electronic-effects:intermediates}{कार्बधनायन} का
स्थायित्व), ब्रोमोमेथेन और 1-ब्रोमोप्रोपेन शायद ही अभिक्रिया करते हैं।
\end{solution}

\begin{solution}{exo:b1:nucleophilic-substitution:9}
$x_{\text{प्रतीपन}} - x_{\text{धारण}} = 0.12$ और $x_{\text{प्रतीपन}} + x_{\text{धारण}} = 1$:
\qty{56}{\percent} प्रतीपित, \qty{44}{\percent} धारित। अधिकांशतः SN1
\omterm{def:b1:nucleophilic-substitution:sn1}{रेसेमीकरण}, प्रतीपन के थोड़े आधिक्य के साथ: जल के पहुँचने पर निकलता हुआ ब्रोमाइड
अब भी अपनी ओर को ढके रहता है।
\end{solution}

\begin{solution}{exo:b1:nucleophilic-substitution:10}
$v = k_1k_2[\mathrm{RY}][\mathrm{Nu}]/(k_{-1}[\mathrm{Y^-}] + k_2[\mathrm{Nu}])$
(\cref{prop:b1:nucleophilic-substitution:sn1-law})। \ce{Y-} मिलाने से
हर बढ़ता है: अधिक \omterm{def:b1:electronic-effects:intermediates}{कार्बधनायन} ग्रहण होने के बजाय \omterm{def:b1:nucleophilic-substitution:substitution}{क्रियाधार} में
लौट जाते हैं, और दर घटती है। नगण्य, जब
$k_2[\mathrm{Nu}] \gg k_{-1}[\mathrm{Y^-}]$ हो, जैसे जब \omterm{def:b1:electronic-effects:nucleophile}{नाभिकरागी}
विलायक ही हो।
\end{solution}

\begin{solution}{exo:b1:nucleophilic-substitution:11}
समान सांद्रताएँ: $-\mathrm{d}c/\mathrm{d}t = kc^2$, $1/c = 1/C_0 + kt$,
$t_{1/2} = 1/(kC_0) = \qty{5.0e5}{s}$ (लगभग छह दिन)। हाइड्रॉक्साइड के पचास गुना
आधिक्य के साथ $[\ce{OH-}] \approx 50C_0$ स्थिर है: $c = C_0
e^{-k't}$, जहाँ $k' = 50kC_0 = \qty{1.0e-4}{s^{-1}}$, $t_{1/2} = \ln 2/k' =
\qty{6.9e3}{s}$, लगभग दो घंटे।
\end{solution}

\begin{solution}{exo:b1:nucleophilic-substitution:12}
ब्रोमाइड को \ce{OH-} को बाहर निकालना पड़ता, जो \omterm{def:b1:predominant-reaction:strong-acid}{प्रबल क्षारक} और खराब \omterm{def:b1:electronic-effects:nucleophile}{निर्गामी
समूह} है: कोई अभिक्रिया नहीं। सांद्र HBr में OH प्रोटॉनित होकर
\ce{-OH2+} बनता है, जिसका \omterm{def:b1:electronic-effects:nucleophile}{निर्गामी समूह} जल है। तृतीयक ऐल्कोहॉल: जल निकलता है,
जिससे \omterm{def:b1:electronic-effects:intermediates}{कार्बधनायन} बनता है, जिसे \ce{Br-} ग्रहण करता है (SN1)। प्राथमिक ऐल्कोहॉल:
\ce{Br-}, \ce{R-OH2+} के कार्बन पर पीछे से आक्रमण करता है जबकि जल निकलता है
(SN2)।
\end{solution}

\begin{solution}{pb:b1:nucleophilic-substitution:1}
\textbf{1.} अभिक्रिया करने वाला हर अणु एक \ce{H+} और एक \ce{Cl-} देता है;
चालकता आयनिक योगदानों का योग है, जो आयनों की मात्रा के
समानुपाती है।
\textbf{2.} $[\mathrm{RCl}]/[\mathrm{RCl}]_0 = (\kappa_\infty -
\kappa)/\kappa_\infty$।
\textbf{3.} $\ln[\mathrm{RCl}] = \ln[\mathrm{RCl}]_0 - kt$: $\ln(\kappa_\infty - \kappa)$ को
$t$ के विरुद्ध आलेखित कीजिए।
\textbf{4.} 0.693, 0.603, 0.513, 0.423, 0.333, 0.153, $-0.117$, $-0.387$।
\textbf{5.} हाँ, पूरे समय प्रवणता $-0.0180$ प्रति मिनट: प्रथम कोटि।
\textbf{6.} जल में कोई भी कोटि: उसकी सांद्रता स्थिर है
(कोटि का ह्रास)।
\textbf{7.} $k = 0.0180/60 = \qty{3.0e-4}{s^{-1}}$।
\textbf{8.} $t_{1/2} = \ln 2/k = \qty{2.3e3}{s}$, \qty{38.5}{min}।
\textbf{9.} $\ln 10/k = \qty{7.7e3}{s}$, \qty{128}{min}।
\textbf{10.} $\ln(\kappa_\infty - \kappa)$: $t = 0$ पर 0.693,
\qty{10}{min} पर 0.109, इत्यादि: प्रवणता $-0.0584$ प्रति मिनट, $k = \qty{9.7e-4}{s^{-1}}$।
\textbf{11.} $9.7/3.0 = 3.2$।
\textbf{12.} \qty{30}{min} पर: $2.000(1 - e^{-0.0584 \times 30}) = 1.653$,
मापा गया 1.654।
\textbf{13.} SN1: तृतीयक \omterm{def:b1:nucleophilic-substitution:substitution}{क्रियाधार}, \omterm{def:b1:intermolecular-forces-solvents:solvent-classes}{प्रोटिक विलायक}, उदासीन \omterm{def:b1:electronic-effects:nucleophile}{नाभिकरागी}।
\textbf{14.} \cref{exo:b1:nucleophilic-substitution:5} की तरह, Cl के साथ।
\textbf{15.} मंद आयनन में केवल \omterm{def:b1:nucleophilic-substitution:substitution}{क्रियाधार} भाग लेता है: $v =
k[\mathrm{RCl}]$।
\textbf{16.} हाइड्रॉक्साइड केवल तीव्र ग्रहण चरण में, दर-निर्धारक आयनन के
बाद, भूमिका निभा सकता है; दर उस पर निर्भर नहीं करती।
\textbf{17.} \omterm{def:b1:electronic-effects:intermediates}{कार्बधनायन} समतलीय और \omterm{def:b1:stereochemistry-in-depth:stereocentre}{अकिरैल} है; जल दोनों
फलकों पर समान रूप से जुड़ता है: रेसेमिक ऐल्कोहॉल।
\textbf{18.} उथले गर्त में स्थित \omterm{def:b1:electronic-effects:intermediates}{कार्बधनायन} तक एक ऊँचा पहला अवरोध (आयनन),
फिर ग्रहण तक एक नीचा अवरोध: पहला चरण
दर-निर्धारक है।
\textbf{19.} वह एक प्राथमिक \omterm{def:b1:electronic-effects:intermediates}{कार्बधनायन} बनाता (बहुत अस्थायी), और जल
तीव्र SN2 के लिए बहुत दुर्बल \omterm{def:b1:electronic-effects:nucleophile}{नाभिकरागी} है।
\textbf{20.} $k = A\,e^{-E_a/RT}$।
\textbf{21.} $E_a = R\ln(k_2/k_1)/(1/T_1 - 1/T_2)$।
\textbf{22.} $E_a = 8.314 \times \ln(9.74/3.00)/(1/298.15 - 1/308.15) =
\qty{9.0e4}{J/mol}$।
\textbf{23.} आयनन के अवरोध की ऊँचाई, जो
\omterm{def:b1:elementary-steps:rds}{दर-निर्धारक चरण} है।
\textbf{24.} $k = 3.0 \times 10^{-4} \exp\bigl(\frac{90\,000}{8.314}
(\frac{1}{298.15} - \frac{1}{318.15})\bigr) = \qty{2.9e-3}{s^{-1}}$।
\textbf{25.} \textbf{$E_a = \qty{90}{kJ/mol}$}।
\end{solution}
